\section{The Longest Word}

The eternity of hell rests, in the end, on one sentence. Everything else in the New Testament is compatible with a long punishment or a terminal one; this sentence is the place where the Church's reading and the alternative readings actually meet. It closes the last discourse in Matthew, after the sheep and the goats have been separated and the sentence pronounced:

\begin{studyframe}
And these shall go into everlasting punishment: but the just, into life everlasting.\endnote{Matthew 25:46, Douay--Rheims (Challoner revision), read 2026-07-25 in this library's tracked verse-text derivative of the Project Gutenberg Challoner text at its exact locus, together with the whole of 25:31--46. Greek: \latin{kai apeleusontai houtoi eis kolasin ai\=onion; hoi de dikaioi eis z\=o\=en ai\=onion}, read the same day in the tracked Robinson--Pierpont Byzantine Textform CSV for Matthew at 25:46. Greek is transliterated throughout this essay because the typeface used here carries no Greek; the transliterations are the author's and the registered edition and locus are given for every one so that a reader can check them.}
\end{studyframe}

The two clauses carry the same adjective. In Greek it is \latin{ai\=onios} on both sides; in the Vulgate tradition behind the Douay--Rheims it is rendered by two English forms of one Latin word. Whatever the adjective means in the second clause it means in the first, and the sentence gives no purchase for splitting it.

That is the argument. It is not a subtle argument and it is not a modern one. Augustine had it in the early fifth century, deployed it against the identical objection, and stated it in terms nobody has improved on:

\begin{studyframe}
Then what a fond fancy is it to suppose that eternal punishment means long-continued punishment, while eternal life means life without end, since Christ in the very same passage spoke of both in similar terms in one and the same sentence, ``These shall go away into eternal punishment, but the righteous into life eternal!'' If both destinies are ``eternal,'' then we must either understand both as long-continued but at last terminating, or both as endless. For they are correlative,---on the one hand, punishment eternal, on the other hand, life eternal. And to say in one and the same sense, life eternal shall be endless, punishment eternal shall come to an end, is the height of absurdity.\endnote{Augustine, \emph{De civitate Dei} XXI.23; Dods 1871 translation; the complete chapter read 2026-07-25 in this library's tracked transcription at artifact lines 17958--18013 and registered there as a passage. The chapter's earlier argument, that the sentence pronounced on men and angels alike cannot be true of the angels and false of the men, is treated below.}
\end{studyframe}

It is worth pausing on how the argument works, because its structure determines what could defeat it. It is not an argument from the meaning of a word. It is an argument from a symmetry, and it is indifferent to the lexical question: it says that whatever the word means, it cannot mean two different things inside one clause of one sentence, and that a reader who terminates the punishment has undertaken to terminate the life. That is why it survives the lexical objections that follow --- and it is also why those objections deserve to be stated at their full strength, since a reader who has met them only as caricature will not understand why serious people have found them serious.

\subsection{The case against, at full strength}

\textbf{First: the adjective does not mean ``endless,'' and the New Testament itself proves it.} The Greek \latin{ai\=onios} is formed from \latin{ai\=on}, an age or epoch, and its native force is ``belonging to the age'' or ``age-long.'' Greek possessed a stronger word for strict endlessness, \latin{aidios}. The Epistle of Jude uses both, one verse apart. Of the fallen angels: God ``hath reserved under darkness in everlasting chains'' --- \latin{desmois aidiois}. Of the cities of the plain, in the next sentence: they ``were made an example, suffering the punishment of eternal fire'' --- \latin{pyros ai\=oniou}.\endnote{Jude 6--7, Douay--Rheims at its locus; Greek from the tracked Robinson--Pierpont CSV for Jude, read 2026-07-25: v.~6 \latin{eis krisin megal\=es h\=emeras desmois aidiois hypo zophon tet\=er\=eken}; v.~7 \latin{prokeintai deigma, pyros ai\=oniou dik\=en hypechousai}.} And the fire that fell on Sodom is not burning. It burned once, consumed the cities, and went out; the Dead Sea is not on fire. If \latin{ai\=onios} meant ``never-ending'' as a description of the burning, Jude's sentence would be false as a matter of geography. The author who had \latin{aidios} available and did not use it here evidently did not mean by \latin{ai\=onios} what the later Latin tradition made it mean.

\textbf{Second: the noun in Matthew 25:46 means corrective chastisement, not retribution.} The word is \latin{kolasis}, from \latin{kolaz\=o}, whose root sense is to prune or check; it is distinguished in the classical vocabulary from \latin{tim\=oria}, vengeance. Aristotle states the distinction in the \emph{Rhetoric} in one line: \latin{diapherei de tim\=oria kai kolasis: h\=e men gar kolasis tou paschontos heneka estin, h\=e de tim\=oria tou poiountos, hina pl\=er\=oth\=ei} --- ``there is a difference between revenge and punishment; the latter is inflicted in the interest of the sufferer, the former in the interest of him who inflicts it, that he may obtain satisfaction.''\endnote{Aristotle, \emph{Rhetoric} I.10.17 (Bekker 1369b12--14). Greek read 2026-07-25 in the Perseus Digital Library presentation of the text edited by W.~D. Ross (Oxford, 1959); English in the same site's presentation of the translation of J.~H. Freese (Loeb, 1926), first published 1926 and out of copyright in the United States by date. Perseus states that its markup is licensed CC~BY-SA~3.0 United States. Freese's rendering of \latin{kolasis} as ``punishment'' and \latin{tim\=oria} as ``revenge'' is a translator's choice and does not settle the Greek.} If Matthew's phrase means ``age-long corrective chastisement'' rather than ``endless retributive torment,'' the sentence is not merely compatible with a purgative reading; it invites one. And Aristotle is not a fringe witness: he is the standard authority on Greek usage for the period in which the Gospels were written.

\textbf{Third: the strongest Pauline text says destruction, not torment.} Paul's only extended statement on the fate of the wicked runs: they ``shall suffer eternal punishment in destruction, from the face of the Lord and from the glory of his power.'' The Greek is \latin{olethron ai\=onion} --- eternal ruin or destruction.\endnote{2 Thessalonians 1:8--9, Douay--Rheims at its locus; Greek from the tracked Robinson--Pierpont CSV for 2 Thessalonians, registered for this essay: \latin{hoitines dik\=en tisousin, olethron ai\=onion apo pros\=opou tou kyriou kai apo t\=es dox\=es t\=es ischyos autou}. The whole of 1:6--10 was read.} Matthew supplies a second: fear him who can ``destroy both soul and body in hell'' --- \latin{apolesai}, to destroy or lose utterly.\endnote{Matthew 10:28, Douay--Rheims at its locus; Greek from the tracked Matthew CSV: \latin{phob\=eth\=ete de mallon ton dynamenon kai t\=en psych\=en kai to s\=oma apolesai en geenn\=ei}.} On the annihilationist reading these texts mean what they say: the wicked are destroyed, the destruction is permanent (and so ``eternal'' in the only sense that matters), and the doctrine of unending conscious torment is a later importation.

\textbf{Fourth: the Old Testament image behind the fire is a rubbish-heap, not a torture chamber.} The single verse that supplies the New Testament's most vivid phrase reads: ``And they shall go out, and see the carcasses of the men that have transgressed against me: their worm shall not die, and their fire shall not be quenched: and they shall be a loathsome sight to all flesh.''\endnote{Isaias 66:24, Douay--Rheims at its locus, read 2026-07-25 in the tracked verse-text derivative. The Vulgate wording behind ``carcasses'' was not separately collated here, and no claim below rests on the Latin. This is the verse quoted verbatim in Mark 9:47 (Vulgate and Douay--Rheims numbering; 9:48 in the Greek and in modern editions).} What is unquenched there is the fire, and what does not die is the worm; the objects on which they work are corpses. The undying worm is the permanence of the disgrace, not the permanence of a sufferer.

\textbf{Fifth: the Apocalypse's eternity attaches to the smoke.} ``And the smoke of their torments shall ascend up for ever and ever'' --- \latin{kai ho kapnos tou basanismou aut\=on eis ai\=onas ai\=on\=on anabainei}.\endnote{Apocalypse 14:11, Douay--Rheims at its locus; Greek from the tracked Revelation CSV registered for this essay. The rest of the verse is given below.} A rising column of smoke is what a burnt city leaves behind; the image is of a monument, not of a duration of pain.

\textbf{Sixth: the received text of the strongest Marcan passage is doubtful.} Augustine's argument in \emph{City of God} XXI.9, quoted in the previous section, leans on the fact that the Lord ``did not shrink from using the same words three times over in one passage'': the refrain ``where their worm dieth not, and the fire is not extinguished'' appears three times in the Douay--Rheims text of Mark 9. In the modern critical text it appears once. The first two occurrences --- verses 44 and 46 in the Vulgate numbering --- are marked as omissions in the Nestle--Aland text and are generally regarded as assimilation to the surviving third.\endnote{Mark 9:42--47 in the Douay--Rheims and Vulgate numbering (9:43--48 in the Greek), read 2026-07-25 in the tracked verse-text derivative; the Greek read the same day in the Robinson--Pierpont CSV for Mark registered for this essay, where verses 44 and 46 each carry an inline apparatus brace recording that the Nestle--Aland text omits them, and verse 48 carries no such note. The Robinson--Pierpont apparatus is a divergence note and not a full critical apparatus; this essay reports the divergence and does not adjudicate the text. Augustine could not have known of it, and no criticism of him is intended by reporting it.} The rhetorical force Augustine appealed to belongs to a text he inherited rather than to the words of the Lord as the best witnesses transmit them.

\textbf{Seventh: another whole body of texts points the other way.} ``And when all things shall be subdued unto him, then the Son also himself shall be subject unto him that put all things under him, that God may be all in all.'' ``For God hath concluded all in unbelief, that he may have mercy on all.'' ``As by the offence of one, unto all men to condemnation: so also by the justice of one, unto all men to justification of life.'' ``And through him to reconcile all things unto himself, making peace through the blood of his cross.'' ``The times of the restitution of all things.'' ``Who will have all men to be saved and to come to the knowledge of the truth.''\endnote{1 Corinthians 15:28; Romans 11:32; Romans 5:18; Colossians 1:20; Acts 3:21; 1 Timothy 2:4 --- Douay--Rheims, each read at its own locus in the tracked verse-text derivative on 2026-07-25, with 1 Cor 15:22--28 and Rom 5:18--19 read in full context. These are the texts on which Origen built, and the last of them is the one Aquinas has to answer in the article on whether the will of God is always fulfilled.} Every one of those is Scripture, and the reading that makes them all express nothing but an unrealized wish is not obviously the natural one.

\subsection{The case for, given the case against}

None of the seven is frivolous. Two of them are, in this essay's judgment, straightforwardly correct as far as they go, and saying so costs the doctrine nothing.

The fourth is correct: the Isaian image \emph{is} of corpses, and the eternity in Isaias is the eternity of the worm and the fire rather than of a conscious sufferer. What follows is not that the New Testament means the same thing but that the New Testament is doing something with the image, which is what quotation of the prophets normally does. The sixth is correct: Augustine's appeal to the threefold repetition is weakened by the text-critical evidence, and a reader who wants to argue from Mark should argue from 9:47 alone, which stands in all witnesses and is a direct quotation of Isaias.

The others do not hold, and it is worth being exact about why, because the reasons are not the ones usually given.

The Jude objection is the sharpest, and its answer is in the sentence itself. Jude does not say that Sodom's fire burns for ever. He says that the cities \latin{prokeintai deigma} --- are set forth as an exhibit, a specimen --- ``suffering the punishment of eternal fire.'' The historical burning is the exhibit; the eternal fire is what it is an exhibit \emph{of}. Read that way the sentence is about the fire of the judgment and not about the geology of the Dead Sea, and the choice of \latin{ai\=onios} over \latin{aidios} carries no weight, because the two words are used of different things: strict endlessness of the angels' chains, and the fire of the age to come for the cities' example. This is the ordinary reading and it is not special pleading; but it is a reading, and a reader should know that it is one.

The \latin{kolasis} objection proves less than it is asked to prove. Aristotle's distinction is real and is stated where the objection says it is stated. But it is Aristotle's technical usage in a handbook of persuasion, offered as a clarification precisely because ordinary speakers ran the two words together; and by the first century the corrective force of \latin{kolasis} had worn thin enough that the New Testament can use the noun for the pain that accompanies fear, where nothing corrective is in view.\endnote{1 John 4:18, ``fear hath pain'' (Douay--Rheims at its locus, read 2026-07-25); the Greek of 1 John was not reached in a registered edition and the claim about its wording is therefore stated here only as far as the Latin and English witnesses carry it: the Vulgate tradition renders the noun by \latin{poena}, and the sentence is not about correction. A reader for whom the lexical question is decisive should consult a Greek concordance directly; this essay does not rest its conclusion on the point.} More decisively, the objection is aimed at the wrong target. Grant \latin{kolasis} its corrective sense in full: the adjective still governs both members of Matthew 25:46, and a corrective chastisement that is \latin{ai\=onios} in the same sense as the life is a chastisement that never finishes correcting --- which is not an improvement.

The \latin{olethros} objection reads three words and stops before the fourth. Paul does not write ``eternal destruction'' and leave it there; he writes eternal destruction \latin{apo pros\=opou tou kyriou}, ``from the face of the Lord and from the glory of his power.'' A thing that has been annihilated is not \emph{from} anything; ``from the face of the Lord'' names a relation, and a relation requires two terms. The phrase describes exclusion, which is exactly the \latin{poena damni} of the previous section, and it is difficult to construe as extinction without deleting the prepositional phrase that gives it its content.

The Apocalypse objection stops one clause early in the same way. The verse continues: ``neither have they rest day nor night, who have adored the beast and his image.'' It is the smoke that ascends and the worshippers who have no rest, and the second clause has a personal subject. Six chapters later the same book says of the devil, the beast, and the false prophet that ``they shall be tormented day and night for ever and ever'' --- \latin{basanisth\=esontai h\=emeras kai nyktos eis tous ai\=onas t\=on ai\=on\=on} --- with a plural verb and personal subjects and no smoke anywhere in the sentence.\endnote{Apocalypse 14:11 and 20:10, Greek from the tracked Revelation CSV, read 2026-07-25 with the surrounding verses. A versification caution: the Douay--Rheims, following the Vulgate, divides these verses differently from the Greek, so that its 20:9 ends ``where both the beast'' and its 20:10 begins ``And the false prophet shall be tormented.'' The English reader working from the Douay--Rheims will find the sentence broken across the verse boundary; the Greek runs continuously.} This is the passage Augustine cites in XXI.23 beside Matthew 25:41, and it is why he says that in Scripture ``everlasting'' and ``for ever'' ``are wont to mean nothing else than endless duration.''

The seventh objection --- the universal texts --- is the one that cannot be disposed of by closer reading, because the closer reading confirms that they say what they appear to say. The reply is of a different kind, and it should be given as what it is: an account of how a canon containing both sets is to be held together, not a demonstration that one set is illusory. Three points carry it. The context of the greatest of them is subjection rather than salvation: 1 Corinthians 15 is about enemies being put under Christ's feet, and being subjected to a king is not the same as being reconciled to him. The scope of the offer is not the same as the tally of its acceptance: ``who will have all men to be saved'' states what God wills for all, and the whole question of the last section of this essay is what that willing amounts to when it meets a refusal. And the texts of universal reconciliation stand in the same books as the texts of final loss, often within a few pages: the Paul of Romans 11:32 is the Paul of 2 Thessalonians 1:9, and a reading that keeps one by discarding the other has not solved the problem but chosen a side of it.

\subsection{What the definitions do to the argument}

For a Catholic reader the exegetical dispute is not open in the way the foregoing might suggest, and pretending otherwise would be its own kind of dishonesty. The Church has determined the sense: the \latin{Quicumque}, the canon of 543, \emph{Firmiter}, \emph{Cantate Domino}, and Paul VI all read the texts one way, and the last of them does so in 1968, after every one of the seven objections had been in print for a century or more. The doctrine of the eternity of hell does not rest on the lexicography of \latin{ai\=onios} and could not, because no doctrine rests on lexicography. It rests on the Church's reading of the whole, exercised as an act of authority, and the exegetical arguments are the reasons that reading gives for itself rather than the ground on which it stands.

But that is a claim about authority, not a claim that the objections are stupid, and the distinction matters practically. A reader who has been told that the alternative readings are the inventions of people who wanted an easier religion will one day encounter Jude 7 and discover that it is not so, and the discovery will cost him more than the objection was worth. It is better to say plainly: the case against is real, parts of it are correct, the whole of it has been available to the Church for sixteen centuries, and she has read the texts otherwise.

\begin{studybox}{Project synthesis --- an observation about the form of the definitions}
Something in the acts themselves suggests that their drafters had noticed exactly where the argument lies. Matthew 25:46 states the two destinies in one sentence under one adjective. The \latin{Quicumque} states them in one sentence under one adjective: \latin{in vitam aeternam\ldots{} in ignem aeternum}. \emph{Firmiter} states them in one clause with the two datives balanced: \latin{illi cum diabolo poenam perpetuam, et isti cum Christo gloriam sempiternam}. Paul VI in 1968 states them in one sentence with the participles balanced. Four acts across fifteen centuries reproduce the grammatical symmetry of the verse rather than paraphrasing it, and each of them could have said simply that the punishment is everlasting without saying it in a construction that couples it to the life. This is an observation about form and it is offered as nothing more; drafters imitate their sources for many reasons, and the pattern is consistent with mere quotation. But it is at least suggestive that the Church's own formulations have preserved the one feature of the text on which the whole weight of the argument falls.
\end{studybox}
